% ============================================================================
% METHODS BRIEF: DDI KNOWLEDGE GRAPH CONSTRUCTION
% Publication Format - Knowledge Graph Focus
% ============================================================================

\documentclass[11pt,a4paper]{article}
\usepackage[utf8]{inputenc}
\usepackage{amsmath,amssymb}
\usepackage{booktabs}
\usepackage{geometry}
\usepackage{graphicx}
\usepackage{hyperref}
\usepackage{float}
\geometry{margin=1in}

\begin{document}

\section*{Methods}

\subsection*{Knowledge Graph Construction}

We constructed a comprehensive drug-drug interaction (DDI) knowledge graph by integrating data from three authoritative biomedical databases: DrugBank v5.1.9, SIDER 4.1, and the Comparative Toxicogenomics Database (CTD). DrugBank served as the primary source, from which we parsed the complete XML database (1.8 GB) to extract drug identifiers, chemical properties, protein targets, biological pathways, drug categories, and pharmacogenomic data including SNP effects and adverse reactions. Drug-side effect associations were obtained from SIDER 4.1 using MedDRA terminology through exact drug name matching (case-insensitive). Drug-disease relationships were integrated from CTD using a two-tier matching strategy: CAS number matching as the primary method followed by exact drug name matching.

Unlike similarity-based approaches that may introduce spurious relationships, our methodology employs exclusively fact-based linkage using exact identifier matching. The matching hierarchy proceeds as follows: DrugBank unique identifiers serve as the primary matching criterion, Chemical Abstracts Service (CAS) registry numbers provide cross-database linking, and exact case-insensitive name matching serves as the tertiary fallback. All relationships include provenance metadata tracking the data source, matching method, and identifier used.

The knowledge graph was implemented using Python 3.9 with NetworkX 3.1 as the primary graph data structure for in-memory analysis and algorithm execution. Custom parsing scripts processed the DrugBank XML using ElementTree, while pandas DataFrames handled tabular data from SIDER and CTD. The graph is persisted as a NetworkX pickle file for efficient loading and as Neo4j-compatible CSV files (nodes and edges with headers) for graph database deployment and Cypher-based querying.

% Knowledge Graph Schema Figure
\begin{figure}[H]
\centering
\includegraphics[width=0.95\textwidth]{figures/kg_schema.png}
\caption{Neo4j-style knowledge graph schema. Node types (Drug, Protein, Pathway, SideEffect, Disease, Category) display entity counts and property schemas. Relationship types show edge counts and attributes. The DDI severity distribution reflects rule-based recalibration: Moderate (80.1\%), Minor (9.3\%), Contraindicated (5.8\%), Major (4.7\%).}
\label{fig:kg_schema}
\end{figure}

% Data Sources Integration Figure
\begin{figure}[H]
\centering
\includegraphics[width=0.9\textwidth]{figures/fig4_data_sources.png}
\caption{Data source integration overview showing the contributions from DrugBank (primary DDI, drug metadata, protein targets, pathways, and pharmacogenomics), SIDER (side effects via MedDRA), and CTD (disease associations) to the comprehensive knowledge graph.}
\label{fig:data_sources}
\end{figure}

\subsection*{DDI Severity Classification}

The original DrugBank severity classification exhibited severe class imbalance with 56.9\% classified as Contraindicated and 0.0\% as Moderate. To address this limitation, we applied a rule-based classifier validated against the DDInter database, achieving 66.4\% exact accuracy with Cohen's $\kappa$ = +0.096 ($p < 0.0001$). The classifier employs empirically-derived keyword weights based on log-likelihood ratios derived from DDInter training data (n=26,014):

\begin{equation}
\text{score}(d) = \sum_{k \in \mathcal{K}} w_k \cdot \mathbb{1}[k \in d]
\end{equation}

where $w_k = \ln(P(k|\text{Major})/P(k|\text{Moderate}))$ represents the log-likelihood ratio for keyword $k$. Severity levels are assigned using percentile-based thresholds: Contraindicated (top 4\%), Major (top 8\%), Minor (bottom 2\%), and Moderate (remaining).

\section*{Results}

\subsection*{Network Analysis}

The DDI knowledge graph exhibits scale-free network properties characteristic of biological interaction networks. Analysis of the drug-drug interaction subgraph (4,314 drug nodes, 379,917 DDI edges) reveals structural patterns with direct implications for clinical drug safety assessment.

\subsubsection*{Global Network Metrics}

The DDI network demonstrates substantial connectivity with a network density of 0.041, indicating that approximately 4\% of all possible drug pairs exhibit documented interactions. Key topological metrics include:

\begin{table}[H]
\centering
\caption{DDI Network Topological Metrics}
\label{tab:network_metrics}
\begin{tabular}{lr}
\toprule
\textbf{Metric} & \textbf{Value} \\
\midrule
Total Drug Nodes & 4,314 \\
Total DDI Edges & 379,917 \\
Network Density & 0.041 \\
Average Degree & 176.1 \\
Median Degree & 74 \\
Standard Deviation & 321.9 \\
Clustering Coefficient & 0.70 \\
Network Diameter & 5 \\
Average Path Length & 2.19 \\
\bottomrule
\end{tabular}
\end{table}

The high clustering coefficient (0.70) indicates that if drug A interacts with both drugs B and C, then B and C are also likely to interact with each other---a pattern reflecting shared metabolic pathways or receptor targets. The network exhibits ``small-world'' properties: a diameter of 5 means any two drugs in the network are connected by at most 5 edges, while the average path length of 2.19 indicates most drug pairs are connected through approximately 2 intermediate drugs. These properties enable rapid propagation of interaction effects across the pharmacological landscape.

\subsubsection*{Degree Distribution and Hub Drugs}

The degree distribution follows a power-law pattern ($P(k) \propto k^{-\gamma}$, $\gamma \approx 0.74$, $R^2 = 0.71$), characteristic of scale-free networks. This indicates the presence of ``hub'' drugs with disproportionately high interaction counts. The top 10 hub drugs by total interaction count are:

\begin{table}[H]
\centering
\caption{Top 10 Hub Drugs by DDI Count}
\label{tab:hub_drugs}
\begin{tabular}{llr}
\toprule
\textbf{Rank} & \textbf{Drug} & \textbf{Total DDIs} \\
\midrule
1 & Quinidine & 2,391 \\
2 & Procaine & 2,360 \\
3 & Lidocaine & 2,104 \\
4 & Verapamil & 2,104 \\
5 & Propranolol & 2,092 \\
6 & Warfarin & 2,073 \\
7 & Isradipine & 2,071 \\
8 & Pindolol & 2,024 \\
9 & Indomethacin & 2,023 \\
10 & Nicardipine & 2,022 \\
\bottomrule
\end{tabular}
\end{table}

These hub drugs span multiple therapeutic classes including antiarrhythmics (Quinidine), local anesthetics (Procaine, Lidocaine), calcium channel blockers (Verapamil, Isradipine, Nicardipine), beta-blockers (Propranolol, Pindolol), anticoagulants (Warfarin), and NSAIDs (Indomethacin), reflecting their broad pharmacokinetic interaction profiles via CYP450 metabolism and protein binding displacement.

\subsubsection*{Severity-Stratified Network Analysis}

Stratifying the DDI network by interaction severity reveals distinct subnetwork structures. Using our rule-based recalibrated severity classifier (66.4\% DDInter accuracy), the distribution shifts from the original DrugBank's severe bias to a more clinically realistic pattern:

\begin{table}[H]
\centering
\caption{DDI Subnetwork Statistics by Severity (Rule-Based Recalibrated)}
\label{tab:severity_networks}
\begin{tabular}{lrrrrr}
\toprule
\textbf{Severity} & \textbf{Edges} & \textbf{Pct} & \textbf{Density} & \textbf{Avg Degree} & \textbf{Components} \\
\midrule
Moderate & 304,399 & 80.1\% & 0.036 & 148.6 & 1 \\
Minor & 35,343 & 9.3\% & 0.016 & 33.8 & 1 \\
Contraindicated & 22,153 & 5.8\% & 0.084 & 61.1 & 1 \\
Major & 18,022 & 4.7\% & 0.082 & 54.2 & 1 \\
\midrule
\textit{Severe (C+M)} & \textit{40,175} & \textit{10.6\%} & --- & --- & --- \\
\bottomrule
\end{tabular}
\end{table}

The Moderate subnetwork dominates (304,399 edges, 80.1\%), spanning 4,098 drugs. The severe interaction subnetwork (Contraindicated + Major = 40,175 edges, 10.6\%) involves 725 drugs for Contraindicated and 665 drugs for Major interactions. All severity subnetworks form single connected components, meaning drugs within each severity class are reachable from every other---there are no isolated ``islands'' of drugs at any severity level. This underscores the interconnected nature of drug interactions across the pharmacological space.

\subsubsection*{Centrality Analysis}

We computed multiple centrality measures to identify drugs with high-risk interaction profiles (Figure~\ref{fig:centrality_heatmap}):

\textbf{Betweenness Centrality:} Drugs with high betweenness serve as ``bridges'' in the interaction network, mediating indirect effects between drug clusters. High-betweenness drugs include local anesthetics (Procaine, Lidocaine) that interact across multiple drug classes, and biologics (Sotatercept, Abciximab) that bridge distinct therapeutic domains.

\textbf{Eigenvector Centrality:} Unlike degree centrality which counts connections, eigenvector centrality weights connections by the importance of neighboring nodes. A drug scores high if it interacts with other highly-connected drugs. In the DDI network, this identifies drugs that interact with multiple hub medications---adding such a drug to a polypharmacy regimen creates cascading interaction risks because it connects to the most interaction-prone drugs in the network. Antiarrhythmics (Quinidine) and cardiovascular agents (Propranolol, Verapamil, Isradipine) dominate the top eigenvector centrality rankings, reflecting their tendency to interact with other high-risk medications.

\textbf{PageRank:} Originally developed by Google to rank web pages, PageRank in the DDI network measures which drugs are most likely to be ``reached'' by random walks through the interaction graph. A drug has high PageRank if many other drugs interact with it, especially if those drugs themselves have many interactions. In clinical terms, high-PageRank drugs are frequent interaction targets---they are the medications most likely to be affected when new drugs are added to a regimen. Local anesthetics (Procaine, Lidocaine, Benzocaine) exhibit the highest PageRank scores, reflecting their extensive interaction profiles with multiple drug classes via sodium channel and protein binding mechanisms.

% Centrality Heatmap
\begin{figure}[H]
\centering
\includegraphics[width=0.45\textwidth]{figures/centrality_heatmap.png}
\caption{Multi-centrality profile heatmap for top 20 drugs (ranked by average centrality). Normalized values (0-1) enable comparison across degree, betweenness, eigenvector, and PageRank metrics. Procaine, Lidocaine, and Quinidine consistently rank highest across centrality measures, confirming their status as critical hub drugs requiring enhanced monitoring in polypharmacy scenarios.}
\label{fig:centrality_heatmap}
\end{figure}

\subsubsection*{Community Detection}

Application of the Louvain community detection algorithm identified 6 distinct drug communities (modularity Q = 0.23). The moderate modularity reflects the highly interconnected nature of drug interactions across pharmacological classes. The largest community (Community 1) contains 1,235 drugs spanning multiple therapeutic classes, followed by Community 2 with 858 mixed pharmacological agents, Community 3 with 766 cardiovascular and alimentary agents, Community 4 with 727 drugs forming a CNS-enriched cluster, and Community 5 with 708 drugs exhibiting CNS/cardiovascular overlap.

The lack of clean pharmacological class separation underscores that DDIs cross therapeutic boundaries, reinforcing the need for comprehensive interaction checking in polypharmacy scenarios. Inter-community edges represent cross-class interactions with potentially unexpected clinical consequences. The community structure is visualized in Figure~\ref{fig:louvain_network}, showing representative hub drugs from each Louvain community with distinct colors indicating community membership.

% Louvain Network
\begin{figure}[H]
\centering
\includegraphics[width=0.85\textwidth]{figures/louvain_network.png}
\caption{DDI network visualization with Louvain community detection. Node colors indicate community membership (6 communities, modularity Q = 0.23). Representative hub drugs from each major community are shown, with node size proportional to degree (number of interactions). Labels identify highest-degree drugs. The network structure reveals how drugs cluster into interaction neighborhoods that span therapeutic classes.}
\label{fig:louvain_network}
\end{figure}

\section*{Data Availability}

The knowledge graph is available in multiple formats:
\begin{itemize}
    \item NetworkX pickle file (\texttt{enriched\_kg.pkl}) with complete node and edge attributes
    \item Neo4j-compatible CSV files for graph database deployment
    \item JSON statistics file with node/edge counts
\end{itemize}

\section*{References}

\begin{thebibliography}{4}
\bibitem{wishart2018drugbank} Wishart DS, et al. DrugBank 5.0. \textit{Nucleic Acids Res}. 2018;46(D1):D1074-D1082.
\bibitem{kuhn2016sider} Kuhn M, et al. SIDER database. \textit{Nucleic Acids Res}. 2016;44(D1):D1075-D1079.
\bibitem{davis2021ctd} Davis AP, et al. CTD: update 2021. \textit{Nucleic Acids Res}. 2021;49(D1):D1138-D1143.
\bibitem{xiong2022ddinter} Xiong G, et al. DDInter database. \textit{Nucleic Acids Res}. 2022;50(D1):D1200-D1207.
\end{thebibliography}

\end{document}
